\section{The Nijenhuis--Richardson bracket}\label{sec.NR}

Unless otherwise specif\/ied every (super)vector space, symmetric product, linear map etc.\ are considered over the base f\/ield $\Q$ of rational numbers, although everything works in the same way over any f\/ield of characteristic~0.
Given a super vector space $V=V^{0}\oplus V^{1}$, we denote by~$V^{\odot n}$ its symmetric powers and by
$D_{n}(V)=\Hom^*(V^{\odot n+1},V)$ the super vector space of multilinear supersymmetric maps on $n+1$ variables
\begin{gather*} f\colon \ \underbrace{V\times\cdots\times V}_{n+1}\to V. \end{gather*}
Recall that supersymmetry means that
\begin{gather*} f(v_0,\ldots,v_i,v_{i+1},\ldots,v_n)=(-1)^{|v_i||v_{i+1}|}f(v_0,\ldots,v_{i+1},v_{i},\ldots,v_n),\end{gather*}
where $|v|=0,1$ denotes the parity of a homogeneous element $v$.

Given $f\in D_n(V)$ and $g\in D_m(V)$, $n,m\ge 0$, their Nijenhuis--Richardson bracket \cite{NijRich67} is def\/ined as
\begin{gather*} [f,g]=f\wedgebar g-(-1)^{|f||g|}g\wedgebar f \in D_{n+m}(V),\end{gather*}
where
\begin{gather*} f\wedgebar g
(v_{0},\ldots, v_{n+m})= \sum_{\sigma\in S(m+1,n)}\epsilon(\sigma)
f\big(g(v_{\sigma(0)},\ldots, v_{\sigma(m)}), v_{\sigma(m+1)},\ldots,
v_{\sigma(m+n)}\big),\end{gather*}
and $S(m+1,n)$ is the set of shuf\/f\/les of type $(m+1,n)$, i.e., the set of permutations $\sigma$ of $0,\ldots,n+m$
such that $\sigma(0)<\cdots<\sigma(m)$ and $\sigma(m+1)<\cdots<\sigma(m+n)$. The Koszul sign~$\epsilon(\sigma)$ is equal to $(-1)^{\alpha}$, where $\alpha$ is the number of pairs $(i,j)$ such that $i<j$, $\sigma(i)>\sigma(j)$ and $|v_i||v_j|=1$.

Observe that when $f\in D_0(V)$ the product $f\wedgebar g$ is the same as the composition product $f\circ g$ and therefore
the Nijenhuis--Richardson bracket reduces to the usual super commutator on~$D_0(V)$. We denote by $D(V)=\prod\limits_{n\ge 0}D_n(V)$; the Nijenhuis--Richardson brackets induces on~$D(V)$ a structure of Lie superalgebra:
\begin{gather*} \left[\sum_{i=0}^{\infty}f_i ,\sum_{j= 0}^{\infty}g_j\right]=\sum_{n= 0}^{\infty} \sum_{i=0}^n [f_i,g_{n-i}],\qquad f_i,g_i\in D_i(V) .\end{gather*}

\begin{Remark}\label{rem.coderivationNR}
Denoting by $\overline{S^c}(V)=\oplus_{n\ge 1} V^{\odot n}$ the reduced symmetric coalgebra generated by~$V$,
the composition with the natural projection $p\colon \overline{S^c}(V)\to V$ gives a morphism of super vector spaces
\begin{gather*} P\colon \ \Hom^*\big(\overline{S^c}(V),\overline{S^c}(V)\big)\to \Hom^*\big(\overline{S^c}(V),V\big)=D(V),\end{gather*}
and it is well known \cite{LadaMarkl} that $P$ induces an isomorphism of Lie superalgebras
\begin{gather*} P\colon \ \Coder^*\big(\overline{S^c}(V)\big)\xrightarrow{\simeq} D(V),\end{gather*}
where the bracket on $\Coder^*(\overline{S^c}(V))$ is the super commutator.
\end{Remark}

Assume now that $A$ is a commutative superalgebra, then there exists a distinguished sequence $\mu_n\in D_n(A)$, $n\ge 0$, corresponding to multiplication in~$A$:
\begin{gather*} \mu_n(a_0,\ldots,a_n)=a_0a_1\cdots a_n .\end{gather*}

We have
\begin{gather}\label{equ.bracketmuenne}
\mu_n\wedgebar \mu_m=\dbinom{n+m+1}{m+1}\mu_{n+m},\qquad
[\mu_n,\mu_m]=(n-m)\dfrac{(n+m+1)!}{(n+1)!(m+1)!}
\mu_{n+m} .\end{gather}
In fact, the binomial coef\/f\/icient in the f\/irst formula is equal to the cardinality of the set of shuf\/f\/les involved in the formula for the product $\wedgebar$.

